\section{Shapley Values for Data Source Attribution}
\label{sec:explainable-ai-shapley}

For MERIT to be considered as a "Glass-box" system, source-level attribution must be provided to provide a fair assessment to a recruiter as to how a candidate was scored.
Part of that required the system being able to capture the contribution of each data source to the final match score.
To achieve this, we focus on Shapley Values, a concept from Cooperative Game Theory that is used to attribute the winnings of a group of players to the contribution of each player.


\subsection*{Cooperative Game Theory in Recruitment}
To relate this concept to MERIT, we consider the final match score as "winnings", while the different data sources are considered the "players".
The main objective is to determine how each player contributed to the winnings.

The Shapley value for a data source $i$ is defined as the average marginal contribution of that source across all possible combinations (coalitions) of sources:
\begin{equation}
    \phi_i(v) = \sum_{S \subseteq N \setminus \{i\}} \frac{|S|! (n - |S| - 1)!}{n!} [v(S \cup \{i\}) - v(S)]
\end{equation}
Where:
\begin{itemize}
    \item $n$ is the total number of data sources.
    \item $S$ is a subset of sources that does not include source $i$.
    \item $v(S)$ is the score generated using only the sources in subset $S$.
\end{itemize}

\subsection*{Source Attribution and Fairness}
By calculating these values, MERIT can explicitly state the contribution of each data source to the final match score.
In other words, it is possible to determine whether a candidate got a high score due to self-reported claims on their CV, or from evidence-based metrics from a source, such as GitHub.
This level of transparency is critical for mitigating algorithmic bias and providing defensible hiring decisions.
By building on the SHAP (Shapley Additive Explanations) framework \cite{lundberg2017unified}, MERIT ensures that its source-level attribution remains mathematically consistent and locally accurate.
